\begin{table}[!htbp]\centering
\caption{Where inpatient psychiatric use goes after closure}
\label{tab:displacement-compact}
\small
\begin{tabular}{lrrr}
\toprule
Inpatient psychiatric admissions & Pre & Post & Change \\
\midrule
To the closing hospital & 596 & 0 & -596 \\
To other hospitals & 430 & 379 & -51 \\
\midrule
Total (any hospital) & 1{,}026 & 379 & -647 \\
\bottomrule
\multicolumn{4}{p{0.86\textwidth}}{\footnotesize Notes: Catchment-level averages of inpatient psychiatric admissions (ICD-10 F00--F99), over event years $-3$ to $-1$ versus $+1$ to $+3$, for the 30 PNASH psychiatric closures with sufficient flow support. Exposed catchments are municipalities sending $\geq 5\%$ of their pre-closure inpatient flow to the closing hospital. The 596 admissions lost at the closing hospital are not picked up by other hospitals: admissions to other hospitals do not rise (they fall by 51), so total inpatient psychiatric admissions fall by about 63\%. The lost volume is not observed in inpatient psychiatric care; the design observes inpatient admissions only, not outpatient (CAPS) use.}\\
\end{tabular}
\end{table}
